\section*{Right-forward closures: published statement and formal source}
A TRS $R$ is a set of rules $\ell\to r$ such that $\ell\notin V$ and $\operatorname{Vars}(\ell)\supseteq\operatorname{Vars}(r)$. A relation $\to\subseteq T\times T$ on a set $T$ is strongly normalizing w.r.t. a set of starting objects $T_0\subseteq T$, written $\operatorname{SN}(\to,T_0)$, if there is no infinite sequence $t_0\to t_1\to t_2\to\cdots$ with $t_0\in T_0$. If there is no restriction on the set of starting objects we just write $\operatorname{SN}(\to)$.
The set of right-hand sides of a TRS $R$ is denoted by $\operatorname{rhs}(R)$, and $R$ is right-linear if all terms in $\operatorname{rhs}(R)$ are linear. In an overlay TRS all critical pairs $(s,t,p)$ satisfy $p=\varepsilon$. A TRS $R$ is locally confluent if all its critical pairs $(s,t,p)\in\operatorname{CP}(R)$ are joinable, i.e., $s\to_R^*\circ{}_R^*\!\leftarrow t$.
\paragraph{Original Definition 1.} The right-forward closure $\operatorname{RFC}(R)$ is the least set that contains $\operatorname{rhs}(R)$ and is closed under narrowing:
\[\operatorname{RFC}(R)=\{t\mid\exists s\in\operatorname{rhs}(R).\ s\rightsquigarrow_R^*t\}=\rightsquigarrow_R^*(\operatorname{rhs}(R)).\]
The complete author definitions of rewriting, narrowing, linearity, variable-renaming and critical pairs are preserved in the native formal theories and accompanying published PDF. The source keeps the fixed infinite-variable-type and renaming setup explicitly.
\begin{authortheorem}
Let $R$ be a right-linear or locally confluent overlay TRS. Then
\[\operatorname{SN}(\to_R)\quad\text{if and only if}\quad\operatorname{SN}(\to_R,\operatorname{RFC}(R)).\]
\end{authortheorem}
The published theorem attributes the result to Dershowitz and Hoot [10]. The present authors supply new complete formal proofs of the two conditions, collected together here as a single target. A new proof of an existing corpus claim would not be counted as a new problem; the directed full-index searches found no existing claim from this work or these theorem identities.
